\documentclass[10pt,a4paper]{amsart}
\usepackage{amsmath,amssymb,amsthm,graphicx}
\usepackage{verbatim}
\newenvironment{reference}{\comment}{\endcomment}
%\newenvironment{reference}{}{}
\newenvironment{slogan}{\comment}{\endcomment}
\newenvironment{history}{\comment}{\endcomment}

% For commutative diagrams we use Xy-pic
\usepackage[all]{xy}

% We use 2cell for 2-commutative diagrams.
\xyoption{2cell}
% Theorem environments.
%
\theoremstyle{plain}
\newtheorem{theorem}[subsection]{Theorem}
\newtheorem{proposition}[subsection]{Proposition}
\newtheorem{lemma}[subsection]{Lemma}

\theoremstyle{definition}
\newtheorem{definition}[subsection]{Definition}
\newtheorem{example}[subsection]{Example}
\newtheorem{exercise}[subsection]{Exercise}
\newtheorem{situation}[subsection]{Situation}

\theoremstyle{remark}
\newtheorem{remark}[subsection]{Remark}
\newtheorem{remarks}[subsection]{Remarks}

\numberwithin{equation}{subsection}

% Macros
%
\def\lim{\mathop{\mathrm{lim}}\nolimits}
\def\colim{\mathop{\mathrm{colim}}\nolimits}
\def\Spec{\mathop{\mathrm{Spec}}}
\def\Hom{\mathop{\mathrm{Hom}}\nolimits}
\def\Ext{\mathop{\mathrm{Ext}}\nolimits}
\def\SheafHom{\mathop{\mathcal{H}\!\mathit{om}}\nolimits}
\def\SheafExt{\mathop{\mathcal{E}\!\mathit{xt}}\nolimits}
\def\Sch{\mathit{Sch}}
\def\Mor{\mathop{\mathrm{Mor}}\nolimits}
\def\Ob{\mathop{\mathrm{Ob}}\nolimits}
\def\Sh{\mathop{\mathit{Sh}}\nolimits}
\def\NL{\mathop{N\!L}\nolimits}
\def\CH{\mathop{\mathrm{CH}}\nolimits}
\def\proetale{{pro\text{-}\acute{e}tale}}
\def\etale{{\acute{e}tale}}
\def\QCoh{\mathit{QCoh}}
\def\Ker{\mathop{\mathrm{Ker}}}
\def\Im{\mathop{\mathrm{Im}}}
\def\Coker{\mathop{\mathrm{Coker}}}
\def\Coim{\mathop{\mathrm{Coim}}}

% Boxtimes
%
\DeclareMathSymbol{\boxtimes}{\mathbin}{AMSa}{"02}

%
% Macros for moduli stacks/spaces
%
\def\QCohstack{\mathcal{QC}\!\mathit{oh}}
\def\Cohstack{\mathcal{C}\!\mathit{oh}}
\def\Spacesstack{\mathcal{S}\!\mathit{paces}}
\def\Quotfunctor{\mathrm{Quot}}
\def\Hilbfunctor{\mathrm{Hilb}}
\def\Curvesstack{\mathcal{C}\!\mathit{urves}}
\def\Polarizedstack{\mathcal{P}\!\mathit{olarized}}
\def\Complexesstack{\mathcal{C}\!\mathit{omplexes}}
% \Pic is the operator that assigns to X its picard group, usage \Pic(X)
% \Picardstack_{X/B} denotes the Picard stack of X over B
% \Picardfunctor_{X/B} denotes the Picard functor of X over B
\def\Pic{\mathop{\mathrm{Pic}}\nolimits}
\def\Picardstack{\mathcal{P}\!\mathit{ic}}
\def\Picardfunctor{\mathrm{Pic}}
\def\Deformationcategory{\mathcal{D}\!\mathit{ef}}
\title{A determinant map for a two term free presentation detects local freeness of the cokernel}
\author{The Stacks Project Authors}
\date{stacks-project:duality.tex:local-snapshot-20261012; GNU FDL 1.2 or later}
\begin{document}
\maketitle
\noindent\textit{Editorial note.} All delivered mathematical text is the author's own, copied from the pinned Stacks Project chapter file duality.tex; no statement or proof step was rewritten, reordered or removed.\par\medskip
\noindent\textit{Editorial note.} Preamble provenance: the pinned chapter file opens with an input of the shared preamble and carries no preamble text of its own, so the project's single Stacks-wide preamble staged at sources/stacks-duality/preamble.tex is excerpted verbatim under the recipe fixed by LANE\_BRIEF 22.1, as three spans, lines 8-13 (verbatim package and the reference/slogan/history environments), lines 14-18 (Xy-pic setup) and lines 154-217 (the native theorem declarations and mathematical macros). Nothing was added, removed or redefined; the wrapper class and the amsmath/amssymb/amsthm/graphicx base packages are supplied by the assembler.\par\medskip
\noindent\textit{Editorial note.} The counted statement and proof contain no cross reference, no citation and no included file, so the delivered body is self contained.\par\medskip
\noindent\textit{Editorial note.} Canonical location of the counted proposition: Stacks tag 0E9Y (sources/stacks-duality/pinned-original-source.tex line 7727).\par\medskip
\noindent\textit{Allowed standard prerequisites (editorial).}\begin{itemize}
\item Exterior powers of modules over a locally ringed space. The conventions on wedge products used in the statement and at proof lines 7743-7746. Fix the objects and the vanishing convention for negative exponents.
\item Rank of a map modulo the maximal ideal and neighbouring maximality. The rank argument at proof lines 7764-7771. Concludes local freeness of the cokernel from the stalk isomorphism.
\end{itemize}\medskip
\begin{lemma}
\label{lemma-determinant}
Let $X$ be a locally ringed space. Let
$$
\mathcal{E}_1 \xrightarrow{\alpha} \mathcal{E}_0 \to \mathcal{F} \to 0
$$
be a short exact sequence of $\mathcal{O}_X$-modules.
Assume $\mathcal{E}_1$ and $\mathcal{E}_0$ are locally
free of ranks $r_1, r_0$. Then there is a canonical map
$$
\wedge^{r_0 - r_1}\mathcal{F}
\longrightarrow
\wedge^{r_1}(\mathcal{E}_1^\vee) \otimes \wedge^{r_0}\mathcal{E}_0
$$
which is an isomorphism on the stalk at $x \in X$
if and only if $\mathcal{F}$ is locally free of rank $r_0 - r_1$
in an open neighbourhood of $x$.
\end{lemma}

\begin{proof}
If $r_1 > r_0$ then $\wedge^{r_0 - r_1}\mathcal{F} = 0$ by convention
and the unique map cannot be an isomorphism. Thus we may assume
$r = r_0 - r_1 \geq 0$. Define the map by the formula
$$
s_1 \wedge \ldots \wedge s_r \mapsto
t_1^\vee \wedge \ldots \wedge t_{r_1}^\vee \otimes
\alpha(t_1) \wedge \ldots \wedge \alpha(t_{r_1}) \wedge
\tilde s_1 \wedge \ldots \wedge \tilde s_r
$$
where $t_1, \ldots, t_{r_1}$ is a local basis for $\mathcal{E}_1$,
correspondingly
$t_1^\vee, \ldots, t_{r_1}^\vee$ is the dual basis for $\mathcal{E}_1^\vee$,
and $s'_i$ is a local lift of $s_i$ to a section of $\mathcal{E}_0$.
We omit the proof that this is well defined.

\medskip\noindent
If $\mathcal{F}$ is locally free of rank $r$, then it is straightforward
to verify that the map is an isomorphism. Conversely, assume the map
is an isomorphism on stalks at $x$. Then $\wedge^r\mathcal{F}_x$
is invertible. This implies that $\mathcal{F}_x$ is generated by
at most $r$ elements. This can only happen if $\alpha$ has rank
$r$ modulo $\mathfrak m_x$, i.e., $\alpha$ has maximal rank modulo
$\mathfrak m_x$. This implies that $\alpha$ has maximal rank
in a neighbourhood of $x$ and hence $\mathcal{F}$ is locally free
of rank $r$ in a neighbourhood as desired.
\end{proof}

\end{document}
